\chapter{Metodologia}
\label{chap:metodologia}

Este capítulo descreve o protocolo de compressão semântica proposto e o desenho experimental usado para avaliá-lo. Trata-se de uma pesquisa aplicada, de natureza experimental e abordagem quantitativa: todos os critérios de qualidade adotados são computáveis por código, e todos os artefatos intermediários são persistidos e versionados para permitir a reprodução integral dos resultados.

\section{Visão Geral do Protocolo}
\label{sec:visao-geral}

O protocolo separa a geração de um modelo BPMN 2.0 em uma etapa probabilística curta e uma etapa determinística verificável. Na fase de construção do conjunto de dados, o fluxo completo é:

\begin{alineascomnumero}
	\item \textbf{Pré-processamento}: a descrição textual bruta do processo é reestruturada por um LLM de fronteira em um formato padronizado em linguagem natural (atores, fluxo principal, decisões e paralelismos);
	\item \textbf{Extração estruturada}: um segundo LLM de fronteira converte o texto padronizado em um JSON canônico de BPMN, contendo elementos, fluxos de sequência e \textit{gateways}, sem informação de leiaute;
	\item \textbf{Compressão}: um conversor determinístico (\textit{json\_to\_dsl}) lineariza o grafo do JSON na DSL proposta, com garantia formal de preservação topológica;
	\item \textbf{Expansão}: um transpilador determinístico (\textit{dsl\_to\_xml}) expande a DSL para XML BPMN 2.0 completo;
	\item \textbf{Validação}: o XML é validado contra o XSD oficial da especificação BPMN 2.0 e comparado topologicamente com o JSON de origem.
\end{alineascomnumero}

% TODO: figura do pipeline (diagrama dos 5 estágios com as tabelas do SQLite entre eles)

Os pares (descrição textual, DSL) resultantes constituem o corpus de treino. Na fase de avaliação, um modelo de linguagem de pequeno porte especializado gera a DSL diretamente do texto, e as etapas 4--5 reconstituem e verificam o XML. O modelo probabilístico nunca é responsável pela sintaxe do documento final — apenas pelo conteúdo semântico da representação intermediária.

\section{Construção do Conjunto de Dados}
\label{sec:dados}

\subsection{Fontes e Curadoria}
\label{subsec:fontes}

Não existe conjunto público em escala que associe descrições textuais de processos a modelos BPMN em formato de intercâmbio. A estratégia adotada foi coletar apenas o lado de entrada (descrições textuais) de fontes heterogêneas e sintetizar o lado de saída pelo pipeline da Seção~\ref{sec:visao-geral}. As fontes utilizadas são apresentadas na Tabela~\ref{tab:fontes-dados}.

\begin{table}[h!]
	\Caption{\label{tab:fontes-dados} Fontes do conjunto de dados e partições}%
	\IBGEtab{}{%
		\begin{tabular}{lrll}
			\toprule
			Fonte & N & Partição & Papel \\
			\midrule \midrule
			GitLab Handbook (seções curadas) & 728 & \textit{sft} & treino supervisionado \\
			GitLab Handbook (seções lineares) & 172 & \textit{grpo} & etapa opcional de reforço \\
			PET \cite{bellan2022pet} & 44 & \textit{sft} & treino supervisionado \\
			PMo \cite{brissard2025pmo} & 53 & \textit{holdout} & avaliação (nunca treino) \\
			Zenodo \cite{mangler2023textual} & 24 & \textit{holdout} & avaliação (nunca treino) \\
			\bottomrule
		\end{tabular}%
	}{%
	\Fonte{Elaborado pelo autor}%
	\Nota{De 1.047 amostras coletadas, 26 foram excluídas por degradação identificada durante a validação topológica, resultando em 1.021 amostras ativas.}%
}
\end{table}

A principal fonte de volume é o \textit{handbook} público da empresa GitLab, um manual operacional aberto com milhares de seções procedurais. As seções foram coletadas por \textit{web scraping} e submetidas a uma curadoria automática que atribui um escore procedural (presença de passos, decisões e atores) e classifica cada seção em quatro categorias: \textit{ideal} (passos, decisões e atores), \textit{good} (passos e decisões ou atores), \textit{linear} (apenas passos) e \textit{marginal} (sem passos, descartada). Seções \textit{ideal} e \textit{good} compõem a partição de treino supervisionado; seções \textit{linear}, por não exigirem estrutura de decisão, são reservadas como \textit{prompts} para a etapa opcional de aprendizado por reforço.

\subsection{Pré-processamento com LLM}
\label{subsec:preprocessamento}

Cada descrição bruta é reestruturada pelo modelo Kimi K2.6 (acessado via Ollama Cloud) em um formato padronizado que explicita nome do processo, atores, fluxo principal numerado, decisões com os dois caminhos (sim/não) e ações paralelas. O \textit{prompt} instrui o modelo a apenas reorganizar e limpar o texto, sem inventar informação, e a separar blocos independentes quando a descrição contém múltiplos processos. Essa etapa normaliza a variação estilística das fontes e aproxima o texto do vocabulário estrutural do BPMN sem ainda impor formato de máquina.

\subsection{Geração do JSON BPMN Canônico}
\label{subsec:json-canonico}

O texto padronizado é convertido pelo modelo DeepSeek V4 (Ollama Cloud) em um JSON canônico que descreve o processo como grafo: nós tipados (eventos de início e fim, tarefas por tipo, \textit{gateways} exclusivos, paralelos e inclusivos), arestas de fluxo de sequência com condições, e raias de responsabilidade. O JSON não contém leiaute nem identificadores de serialização — apenas a estrutura lógica do processo. Esse artefato cumpre dois papéis: é a fonte da qual a DSL é derivada deterministicamente e é a referência contra a qual a equivalência topológica do XML final é verificada (Seção~\ref{sec:avaliacao}).

Cabe registrar uma limitação assumida: o "gabarito" estrutural de cada amostra é a interpretação do processo feita por um LLM de fronteira, não uma anotação humana. Essa escolha é o que torna viável a escala do corpus, e seu risco é mitigado pela avaliação final contra o conjunto PMo, cujos modelos de referência são curados por especialistas.

\subsection{Partições e Prevenção de Contaminação}
\label{subsec:particoes}

As amostras são particionadas em \textit{sft} (772), \textit{grpo} (172) e \textit{holdout} (77). Duas regras protegem a integridade da avaliação:

\begin{alineas}
	\item o conjunto PMo é exclusivamente de avaliação: nenhum artefato derivado dele participa de qualquer etapa de treino;
	\item as 24 descrições da fonte Zenodo \cite{mangler2023textual} provêm da mesma origem dos itens 25--48 do PMo (com pré-processamento distinto). Para evitar contaminação indireta do \textit{holdout} — treinar com versões alternativas dos mesmos processos usados na avaliação — essa fonte foi igualmente excluída do treino e alocada à partição \textit{holdout}, onde suas múltiplas modelagens de referência com pontuação de especialistas apoiam a análise de ambiguidade.
\end{alineas}

\section{A BPMN-DSL}
\label{sec:dsl}

A representação intermediária é uma DSL projetada neste trabalho, uma vez que não há linguagem consolidada para descrição textual compacta de BPMN. Seu desenho segue princípios explícitos: cada conceito BPMN tem exatamente uma forma sintática; nomes são legibilidade e identificadores são referência (elementos referenciados recebem sufixo \texttt{\#id} estável); raias são escopos que envolvem o fluxo (bloco \texttt{@lane}), eliminando a dupla declaração de atividades; ramos vazios de decisão são explícitos (\texttt{()}); e propriedades adicionais usam sintaxe única de pares chave-valor. A gramática é formalizada em EBNF e implementada com o \textit{parser} Lark\footnote{\url{https://github.com/lark-parser/lark}}.

A DSL cobre os elementos essenciais da BPMN 2.0 — eventos (com gatilhos de mensagem, temporizador, erro, sinal e escalonamento), oito tipos de tarefa, \textit{gateways} exclusivos, paralelos, inclusivos e baseados em eventos, subprocessos, raias, \textit{pools}, colaborações com fluxos de mensagem e anotações — e mapeia cada construção para o elemento XML correspondente. Conceitos avançados (eventos de borda, compensação, multi-instância) ficam explicitamente fora do escopo desta versão. O Código~\ref{lst:dsl-exemplo} ilustra um processo com decisão e repetição.

\begin{lstlisting}[caption={Exemplo de processo na BPMN-DSL: decisão exclusiva com laço de repetição via referência},label={lst:dsl-exemplo}]
process "Processamento com Retry" {
  start
  -> service "Buscar dados" #fetch
  -> service "Validar"
  -> xor "Valido?" {
    ["sim"] -> service "Processar" -> end
    ["nao"] -> xor "Tentar novamente?" {
                 ["sim"] -> #fetch
                 ["nao"] -> end:error "Invalido"
               }
  }
}
\end{lstlisting}

A meta de projeto é uma razão de compressão de pelo menos 85\% em relação ao XML BPMN canônico correspondente; a medição empírica dessa razão integra o protocolo de avaliação (Seção~\ref{subsec:metricas-modelo}).

\section{Transpilação Determinística}
\label{sec:transpilacao}

\subsection{JSON para DSL}
\label{subsec:json-to-dsl}

O conversor \textit{json\_to\_dsl} lineariza o grafo do JSON canônico em texto DSL. O componente central é o linearizador, que deve satisfazer um único invariante:

\begin{citacao}
Todo nó alcançável do grafo é emitido exatamente uma vez; qualquer reencontro posterior com um nó já emitido é serializado como referência (\texttt{\#id}), nunca como nova emissão.
\end{citacao}

Esse invariante é o que garante que estruturas não estritamente aninhadas — caudas compartilhadas entre ramos, laços com arestas de retorno e convergências múltiplas — sejam representáveis sem duplicação de nós nem perda de lógica. A implementação final sustenta o invariante com três mecanismos: (i) um conjunto único de nós emitidos, compartilhado por referência em toda a travessia recursiva, de modo que cada ramo enxerga o que os demais já emitiram; (ii) detecção determinística de pontos de convergência por pós-dominador imediato, em substituição a heurísticas de busca de junção; e (iii) uma rede de segurança que, ao final da travessia, compara o conjunto de nós alcançáveis a partir do início com o conjunto efetivamente emitido — qualquer diferença é emitida em um apêndice conectado por referências e sinalizada com um aviso rastreável, transformando qualquer perda silenciosa em falha visível.

O conversor foi desenvolvido iterativamente em oito versões, cada uma executada contra o corpus completo e registrada em banco com estado, avisos e erros por amostra, o que permite auditar regressões entre versões. A versão final converte 100\% das 1.021 amostras ativas com preservação topológica verificada.

\subsection{DSL para XML BPMN 2.0}
\label{subsec:dsl-to-xml}

O transpilador \textit{dsl\_to\_xml} expande a DSL para o XML BPMN 2.0 explícito, sem reinterpretar o processo: cada conexão gera um \texttt{sequenceFlow}; blocos de \textit{gateway} geram os nós de divisão e, quando há continuação, de convergência; escopos \texttt{@lane} materializam \texttt{laneSet}, \texttt{lane} e \texttt{flowNodeRef}; colaborações geram participantes e \texttt{messageFlow}. Um contexto global de alocação garante unicidade de identificadores em documentos com múltiplos processos. O XML resultante é validado com a biblioteca lxml contra o XSD oficial da OMG versionado no repositório; a validação por esquema é tratada como critério binário obrigatório de aceitação, e não como métrica.

\section{Protocolo de Avaliação}
\label{sec:avaliacao}

A avaliação é organizada em eixos independentes, para que se possa atribuir cada tipo de falha ao componente responsável.

\subsection{Eixo 1: Validade Sintática}
\label{subsec:eixo1}

Todo XML produzido — na construção do corpus ou na inferência do modelo — deve ser bem formado e válido contra o XSD BPMN 2.0. Para a saída do modelo especializado, mede-se adicionalmente a taxa de DSL sintaticamente válida (aceita pela gramática Lark).

\subsection{Eixo 2: Equivalência Topológica}
\label{subsec:eixo2}

A validade por esquema não garante que o XML represente o mesmo processo do JSON de origem. O segundo eixo compara os dois grafos por: contagem de nós por tipo (eventos, tarefas, \textit{gateways}); conjunto de arestas de fluxo de sequência, do qual derivam precisão, revocação e medida F1; e um indicador binário de equivalência exata de arestas. Esses valores são computados e persistidos para cada amostra do corpus, e o subconjunto com equivalência exata é o utilizado como dados de treino supervisionado de alta confiança. Os resultados dessa verificação sobre o corpus completo são reportados no capítulo de resultados.

\subsection{Métricas sobre a Saída do Modelo}
\label{subsec:metricas-modelo}

A saída do modelo especializado é avaliada contra os modelos de referência do PMo sob quatro dimensões:

\begin{alineas}
	\item \textbf{validade}: taxa de DSL aceita pela gramática e de XML válido contra o XSD;
	\item \textbf{fidelidade semântica}: precisão, revocação e F1 sobre elementos de processo extraídos (atividades, atores, \textit{gateways} e arestas), comparados aos elementos da referência;
	\item \textbf{estrutura}: distância de edição de grafos \cite{sanfeliu1983distance}, em formulação aproximada \cite{gao2010graph}, entre o grafo gerado e o de referência;
	\item \textbf{compressão}: razão de compressão de tokens (TCR) entre a representação gerada e o XML equivalente, medida com o tokenizador do modelo avaliado. A TCR é reportada em três comparações — contra o XML completo, contra o XML sem a seção de diagrama (BPMNDI, que é gerada deterministicamente) e contra o JSON canônico — para que o ganho atribuído à DSL não seja inflado por conteúdo que nenhuma abordagem exigiria do modelo.
\end{alineas}

Como uma mesma descrição textual admite múltiplas modelagens válidas \cite{vanderaa2018challenges}, a porção do \textit{holdout} derivada de \citeonline{mangler2023textual} — que oferece múltiplos modelos de referência com pontuação de especialistas por descrição — é usada para analisar a sensibilidade das métricas à ambiguidade de modelagem.

\subsection{\textit{Baselines}}
\label{subsec:baselines}

O protocolo é comparado sob as mesmas métricas em três configurações: (A) LLM de fronteira gerando XML BPMN diretamente a partir do texto; (B) LLM de fronteira gerando a DSL proposta, com expansão pelo transpilador; e (C) o modelo pequeno especializado gerando a DSL. O contraste A$\times$B isola o efeito da representação intermediária; o contraste B$\times$C isola o efeito da especialização do modelo pequeno frente ao modelo de fronteira.

\section{Especialização do Modelo}
\label{sec:treinamento}

\subsection{Ajuste Supervisionado}
\label{subsec:sft-metodo}

O modelo base é o Qwen2.5-Coder-7B \cite{hui2024qwen25coder}, um modelo aberto de 7 bilhões de parâmetros orientado a código — perfil adequado a uma tarefa de emissão de linguagem formal. A especialização usa ajuste supervisionado com adaptação de baixo posto (LoRA) \cite{hu2022lora} sobre pesos quantizados em 4 \textit{bits} (QLoRA) \cite{dettmers2023qlora}, viabilizando o treino em GPU de 12~GB na fase exploratória, com o treinamento final executado em infraestrutura de nuvem. Os exemplos de treino são os pares (descrição textual bruta, DSL) da partição \textit{sft} com equivalência topológica exata verificada pelo Eixo 2 — o texto de entrada é a descrição original, e não o texto pré-processado, de modo que o modelo especializado absorva também a etapa de estruturação e o sistema final dependa de uma única inferência.

\subsection{Etapa Opcional: Aprendizado por Reforço com Recompensa Verificável}
\label{subsec:grpo-metodo}

Como extensão opcional ao escopo deste trabalho, prevê-se uma etapa de \textit{Group Relative Policy Optimization} \cite{shao2024deepseekmath} sobre o modelo ajustado, usando como \textit{prompts} as seções da partição \textit{grpo} e as amostras sem equivalência exata (que não exigem DSL de referência). A recompensa composta é definida como

\begin{equation}
	r \;=\; 0{,}35\, r_{\text{sint}} \;+\; 0{,}30\, r_{\text{sem}} \;+\; 0{,}25\, r_{\text{topo}} \;+\; 0{,}10\, r_{\text{comp}},
\end{equation}

\noindent em que $r_{\text{sint}}$ (validade da DSL e do XML), $r_{\text{topo}}$ (coerência topológica) e $r_{\text{comp}}$ (economia de tokens) são integralmente verificáveis por código, enquanto $r_{\text{sem}}$ (adequação semântica ao texto) requer sinal aproximado e é o componente metodologicamente mais frágil. Por essa razão, e por restrição de cronograma, essa etapa é tratada como trabalho futuro caso não seja concluída no prazo do TCC.

\section{Infraestrutura e Reprodutibilidade}
\label{sec:reprodutibilidade}

Todo o pipeline é implementado em Python 3.12, com dependências gerenciadas por \textit{uv}, testes automatizados com \textit{pytest} e análise estática com \textit{ruff}. Os artefatos de todas as etapas — texto bruto, texto pré-processado, JSON canônico, DSL, XML e métricas de avaliação — são persistidos em um banco SQLite único, com cada execução de conversor registrada com número de versão, estado e diagnósticos, o que permite auditar e reproduzir qualquer resultado intermediário. A fase exploratória utiliza uma estação com GPU NVIDIA RTX 2080~Ti (12~GB), e o treinamento final, aceleradores de nuvem. Ao término do trabalho, o conjunto de dados e o modelo especializado serão disponibilizados publicamente na plataforma Hugging Face, juntamente com o código-fonte do protocolo.
